\chapter{Restricciones y Metodología}
\label{cap:restricciones_metodologia}

El desarrollo de Luxe Core AI como TFG impone un marco de trabajo delimitado por restricciones técnicas, temporales y económicas que han moldeado las decisiones de diseño. Este capítulo las detalla y explica los \textit{trade-offs} asumidos.

\section{Restricción Temporal}

El límite de un curso académico es la restricción más inflexible del proyecto. Define de forma estricta el alcance del producto mínimo viable (MVP). La metodología adoptada prioriza la iteración rápida y las funcionalidades troncales operativas frente a la sobreingeniería: es mejor tener un sistema que funcione con pocos dispositivos que uno diseñado para mil que nunca llegue a ejecutarse.

\section{Restricción Económica}

El proyecto ha operado con un presupuesto de hardware deliberadamente reducido. El autor es estudiante sin ingresos, lo que ha llevado a seleccionar los dispositivos más económicos disponibles (como el ventilador FanLamp F8808, el modelo más barato de AliExpress con conectividad) y a reutilizar hardware existente (el ESP32 DevKit V1, originalmente adquirido para el puente HVAC, se empleó también como radio BLE programable para la ingeniería inversa del ventilador).

Se ha descartado cualquier compra de hardware auxiliar no estrictamente imprescindible: nada de dongles analizadores BLE, mandos universales o \textit{sniffers} de protocolos. Esta restricción no es una limitación del proyecto, sino una característica deliberada. Obliga a desarrollar competencias de ingeniería inversa que son transferibles a cualquier dispositivo IoT futuro, y demuestra que la domótica inteligente no tiene por qué ser un lujo tecnológico.

\section{Evolución de la Infraestructura}

El desarrollo de modelos de lenguaje locales requiere capacidad de cómputo que un portátil de estudiante no siempre puede proporcionar. Durante la primera fase del proyecto, se utilizó la infraestructura de la Universidad de Córdoba (servidores con GPU RTX, accesibles por VPN y SSH) para prototipar la arquitectura de agentes. Esta fase fue útil para validar el enfoque, pero dependía de una conexión externa y de recursos que el autor no controlaba.

En la segunda fase, el sistema migró a un servidor doméstico dedicado (torre Ryzen APU, 16 GB RAM). Esta migración no fue un retroceso, sino la materialización del objetivo central del proyecto: eliminar cualquier dependencia externa. Hoy, el sistema funciona al 100\% en local, con modelos de lenguaje ejecutándose en Ollama sobre el propio servidor, sin necesidad de GPU dedicada ni de conexión a Internet.

\section{Control de Versiones y Propiedad Intelectual}

El código fuente de Luxe Core AI no es solo un entregable académico: representa la base sobre la que podría construirse una futura iniciativa empresarial. Por eso se aloja en un repositorio privado de GitHub con control de versiones Git. Esta infraestructura no actúa únicamente como mecanismo de \textit{backup} o \textit{rollback}, sino que audita cada contribución y restringe el acceso público al núcleo del proyecto. La trazabilidad que ofrece Git permite, además, vincular cada decisión de diseño documentada en esta memoria con el código que la implementa.

\section{Paradigma Políglota}

La naturaleza dual de la domótica con IA (tiempo real para el control hardware, flexibilidad cognitiva para la toma de decisiones) hace inviable un sistema monolítico en un solo lenguaje. Se ha adoptado una arquitectura de tres capas con lenguajes especializados:

\begin{itemize}
    \item \textbf{Python 3.12} para los \textit{drivers} de dispositivos (Tuya, BLE, FanLamp) y el enrutador de modelos. Python ofrece un ecosistema maduro de librerías (\texttt{tinytuya}~\cite{tinytuya}, \texttt{bleak}~\cite{bleak}) y facilidad de prototipado.
    \item \textbf{Node.js 24} para el orquestador OpenClaw~\cite{openclaw}, que proporciona los canales de usuario (Telegram, \textit{dashboard} web) y la interfaz de agentes de IA.
    \item \textbf{C++20} para el firmware del ESP32, donde se necesita control determinista de memoria y baja latencia en la comunicación con el bus LIN. \textbf{[Pendiente de implementar.]}
\end{itemize}

Esta separación permite elegir el lenguaje óptimo para cada capa sin acoplar los tiempos de ejecución ni forzar compromisos de rendimiento. Herramientas de asistencia de código como OpenCode o Claude Code pueden ejecutarse directamente en el servidor a través de la sesión SSH, beneficiándose de su capacidad de cómputo sin la latencia que implicaría un IDE remoto.

\subsection{Evolución de la Selección de Modelos}

La configuración actual de modelos (1.5B como clasificador ligero, 7B como razonador, con \texttt{bge-m3}~\cite{bge_m3} para \textit{embeddings}) no fue la primera opción. Inicialmente se probó una configuración con un modelo 14B como razonador principal y el 7B como modelo ligero, manteniendo \texttt{bge-m3} para \textit{embeddings}. La idea era que un modelo más grande proporcionaría mejor comprensión contextual.

Sin embargo, las pruebas en el hardware real (Ryzen 5, 16~GB RAM) mostraron que el modelo 14B introducía latencias superiores a un minuto para consultas simples, incluso con cuantización Q4\_K\_M. El consumo de RAM superaba los 12~GB durante la inferencia, dejando apenas margen para el resto del sistema. Se decidió entonces invertir los roles: el 7B pasó a ser el razonador, y se incorporó \texttt{qwen2.5-coder:1.5b} como clasificador ligero para el Tier~1. El modelo \texttt{bge-m3} se mantuvo desde el inicio por su eficiencia (apenas 500~MB en RAM) y su capacidad multilingüe, pues el autor trabaja tanto en español como en inglés.

Esta evolución ilustra una lección importante del proyecto: en el contexto de la \textit{Edge AI}, el equilibrio entre capacidad del modelo y recursos disponibles es tan determinante como la precisión teórica.

\subsection{Base de Datos Vectorial: Decisión de no Incorporar}

Durante el diseño se evaluó la posibilidad de incorporar una base de datos vectorial para mejorar la recuperación semántica de comandos y contextos. Se consideraron tres opciones: \texttt{hnswlib} (librería ligera sin servidor), \texttt{ChromaDB} (base de datos vectorial embebida) y \texttt{Qdrant} (motor especializado con alta eficiencia).

Se decidió no incorporar ninguna en esta fase. El volumen de datos del ecosistema actual es reducido (15 intenciones clasificadas, menos de 200 entradas en caché LRU), y una base de datos vectorial añadiría una capa de complejidad y consumo de RAM no justificados por la ganancia de rendimiento. Los \textit{embeddings} de \texttt{bge-m3} precalculados al arranque y la caché LRU existente cubren las necesidades actuales. Queda como línea futura para cuando el sistema escale a decenas de dispositivos y cientos de comandos.

\section{Aislamiento de Red: La Arquitectura Air-Gapped}

El compromiso con la soberanía del dato exigía un paso adicional. Durante las primeras pruebas se observó que los dispositivos comerciales (Xiaomi, Tuya) intentaban abrir canales salientes hacia sus respectivas nubes en cuanto se conectaban a la red doméstica. La presencia de esa telemetría no controlada erosionaba la promesa fundacional del ecosistema.

La solución fue desplegar una arquitectura de red \textbf{100\% air-gapped}: un segmento \texttt{192.168.1.0/24} sobre un router TP-Link TL-MR6400 sin tarjeta SIM (la ranura está físicamente rota) y sin ningún tipo de enlace WAN. Este router actúa exclusivamente como \textit{switch} Ethernet y punto de acceso WLAN. Todos los dispositivos del ecosistema (servidor, portátil de desarrollo, enchufe Tuya, sensores BLE) conviven en este segmento aislado. El aislamiento físico es la única garantía verificable de que ningún paquete generado por los dispositivos comerciales pueda alcanzar Internet.

\subsection{Hotspot Spoofing para Dispositivos Tuya}

El aislamiento físico introduce una fricción operativa: los dispositivos Tuya necesitan conectarse a Internet durante el emparejamiento inicial para validar el registro contra la nube del fabricante. Para resolverlo sin renunciar al air-gapped, se aplica la técnica de \textit{Hotspot Spoofing}:

\begin{enumerate}
    \item Se levanta temporalmente un punto de acceso móvil (teléfono Android) que replica el SSID y la contraseña del router aislado.
    \item Se completa el emparejamiento del dispositivo Tuya contra la nube del fabricante.
    \item Durante el emparejamiento, se extrae la \textit{Local Key} del dispositivo mediante \texttt{tinytuya wizard}.
    \item Se apaga el \textit{hotspot}. El dispositivo, al encontrar el SSID original (el router TP-Link), se conecta al segmento air-gapped.
    \item A partir de ese momento, todas las órdenes viajan directamente a la IP local del dispositivo (\texttt{192.168.1.100}) mediante el protocolo 3.5 de \texttt{tinytuya}, sin volver a depender de la infraestructura de Tuya.
\end{enumerate}

La \textit{Local Key} queda almacenada exclusivamente en el servidor local y nunca se transmite a terceros.

\subsection{Lectura de Sensores Xiaomi LYWSD03MMC}

Los sensores de temperatura Xiaomi presentaban un problema similar: el firmware de fábrica cifra los anuncios BLE con una \textit{bind key} vinculada a la cuenta del fabricante, lo que reintroduce dependencia de la nube. La ruta prevista era flashear el firmware alternativo ATC MiThermometer~\cite{pvvx_atc_mithermometer}, pero durante las pruebas de campo se descubrió una alternativa mejor.

Al conectar al sensor mediante \texttt{bleak}~\cite{bleak} y obtener un volcado de sus servicios GATT, se descubrió que la característica \texttt{ebe0ccc1} del servicio propietario de Xiaomi es legible \textbf{sin autenticación}. Devuelve 5 bytes con temperatura ($\times$~100), humedad relativa y voltaje de batería:

\begin{itemize}
    \item Bytes 0--1: temperatura en \texttt{int16} LE, con signo. Ejemplo: \texttt{0x09cb} = 2507 $\rightarrow$ 25.07$^\circ$C.
    \item Byte 2: humedad relativa (\%). Ejemplo: \texttt{0x3b} = 59\%.
    \item Bytes 3--4: voltaje de batería en mV (\texttt{uint16} LE). Ejemplo: \texttt{0x0bef} = 3055 mV.
\end{itemize}

Esta vía se adoptó como método principal de integración. El flasheo ATC se mantiene como alternativa para escenarios futuros que requieran lectura pasiva de múltiples sensores simultáneos. La decisión preserva la integridad del hardware (sin riesgo de \textit{brick}) y elimina cualquier dependencia de la nube de Xiaomi. El \textit{trade-off} es que la conexión GATT consume unos 12 segundos por ciclo de lectura, durante los cuales el sensor no puede ser leído por otro dispositivo. En un escenario con un único punto de recolección, esta limitación es aceptable.

\section{Selección del Sistema Operativo y el Servidor}

La inferencia de modelos de lenguaje mediante Ollama demanda recursos que no pueden coexistir de forma estable con un entorno de desarrollo interactivo en el mismo hardware. Por eso se adoptó una arquitectura de dos nodos físicos:

\begin{itemize}
    \item \textbf{Servidor dedicado:} torre con APU Ryzen y 16 GB de RAM, ejecutando \textbf{Ubuntu Server 24.04 LTS} en modo \textit{headless} (sin interfaz gráfica). Alberga el \textit{runtime} completo del ecosistema: Ollama, OpenClaw, los \textit{drivers} Python de dispositivos y la base de datos SQLite. La ausencia de escritorio libera entre 500 y 800 MB de RAM que se destinan a los modelos de lenguaje.
    \item \textbf{Portátil de desarrollo:} conserva las herramientas de ingeniería (compilación de firmware, redacción \LaTeX, Git) y se conecta al servidor exclusivamente por SSH con llaves RSA. Esta separación garantiza que los picos de inferencia no compitan con las tareas interactivas del desarrollador, y viceversa.
\end{itemize}

\subsection{Por qué Ubuntu Server 24.04 LTS}

Se eligió Ubuntu Server por su soporte nativo para todas las dependencias del proyecto (Python 3.12+, Node.js 24, \texttt{bluez} para BLE, \texttt{systemd}) y su ciclo de soporte extendido de 5 años. Dentro de Ubuntu, se optó por la versión 24.04 LTS frente a la recién publicada 26.04 LTS por tres razones:

\begin{enumerate}
    \item \textbf{Madurez del ecosistema:} Ollama, Python, \texttt{bluez} y \texttt{systemd} llevan más de dos años verificados en 24.04. La versión 26.04 introduce versiones renovadas cuyo comportamiento en entornos air-gapped no está suficientemente probado.
    \item \textbf{Disponibilidad de \textit{wheels} binarios:} las librerías Python con dependencias nativas (\texttt{bleak}, \texttt{tinytuya}) tienen paquetes precompilados estables para 24.04. En 26.04 podrían requerir compilación desde fuentes, algo inviable sin acceso a Internet.
    \item \textbf{Estabilidad:} los primeros meses de cualquier LTS concentran la mayor densidad de errores críticos. Operar sobre 24.04 reduce drásticamente el riesgo de encontrar un fallo en un componente del sistema que, en un entorno air-gapped, requeriría intervención física.
\end{enumerate}

Se consideraron otras opciones que fueron descartadas. \textbf{Linux Mint}, incluso en su variante ligera XFCE, consume entre 400 y 700 MB de RAM en reposo debido a los servicios gráficos. Es memoria que en un servidor \textit{headless} se destina íntegramente a los modelos de lenguaje. Las distribuciones \textit{rolling release} como \textbf{CachyOS} o \textbf{Arch Linux} presentan un riesgo adicional en un entorno aislado: una actualización del kernel o de \texttt{bluez} podría romper la compatibilidad con el adaptador Bluetooth, y al no tener acceso a Internet para resolverlo, el sistema quedaría inoperativo hasta una intervención física. Ambos enfoques resultaron incompatibles con los requisitos de estabilidad del proyecto.

\subsection{Herramientas de Gestión}

Para la administración cotidiana del servidor se despliega \textbf{Cockpit} (puerto 9090), una interfaz web ligera que permite monitorizar servicios, CPU y RAM sin necesidad de acceso físico. El desarrollo se realiza mediante \textbf{SSH con autenticación por par de llaves RSA} desde el portátil, lo que permite despliegue remoto de código (\texttt{git pull}) y administración de servicios sin interfaz gráfica.

\section{Despliegue Inicial: Incidencias y Soluciones}

La puesta en marcha del servidor en la red air-gapped enfrentó dos obstáculos que ilustran la clase de problemas que surgen al trabajar con hardware real.

\subsection{El Puerto LAN/WAN Híbrido}

Al conectar el servidor por Ethernet al router TP-Link, el portátil (conectado por Wi-Fi al mismo router) devolvía un error de enrutamiento al intentar acceder por SSH. El cable se había conectado por error al puerto híbrido LAN/WAN del router. Este puerto está diseñado para conectarse a un módem externo, y la lógica interna del \textit{firmware} del TP-Link, al no detectar una conexión WAN válida, mantiene la interfaz en un estado de aislamiento eléctrico que impide el intercambio de tramas Ethernet a nivel de Capa~2 con el resto del \textit{switch} interno.

La solución fue migrar el cable a un puerto LAN dedicado. El \textit{switch} interno restableció inmediatamente el intercambio de tramas, y el servidor apareció en la red con su IP estática \texttt{192.168.1.50}.

\subsection{Pasarela de Red de Contingencia}

La instalación mínima de Ubuntu Server 24.04 LTS carece de herramientas de edición de texto y del demonio \texttt{wpasupplicant} para gestionar la antena Wi-Fi. Sin conectividad a Internet en el segmento air-gapped, el servidor no podía instalar las dependencias del proyecto.

Se implementó una solución temporal utilizando el portátil de desarrollo como \textbf{pasarela NAT} entre el segmento air-gapped y una conexión celular externa (anclaje USB desde un smartphone). La configuración se realizó en dos frentes.

\textbf{En el servidor (Netplan):} se estableció la IP estática y se definió al portátil como \textit{default gateway}:

\begin{verbatim}
network:
  version: 2
  renderer: networkd
  ethernets:
    enp3s0:
      addresses: [192.168.1.50/24]
      routes:
        - to: default
          via: 192.168.1.101
      nameservers:
        addresses: [8.8.8.8, 8.8.4.4]
\end{verbatim}

\textbf{En el portátil:} se ejecutaron los comandos para habilitar el reenvío de paquetes y el enmascaramiento NAT:

\begin{verbatim}
sudo sysctl -w net.ipv4.ip_forward=1
sudo iptables -t nat -A POSTROUTING \
  -o enp0s20f0u2u2 -j MASQUERADE
sudo iptables -P FORWARD ACCEPT
\end{verbatim}

La reconfiguración de la cadena \texttt{FORWARD} fue necesaria porque la política por defecto de \texttt{iptables} en las estaciones de trabajo es \texttt{DROP}: cualquier paquete que intente cursarse entre interfaces distintas del mismo \textit{host} es descartado. Durante la ventana de contingencia, reducir esta protección es seguro porque el segmento air-gapped no tiene salida a Internet en operación normal.

Con esta pasarela, el servidor pudo ejecutar \texttt{apt update \&\& apt dist-upgrade} e instalar todo el \textit{stack} de desarrollo: Python 3.12+, Node.js 24, Git, Cockpit y \texttt{bluez}/\texttt{btmgmt}. Una vez completada la instalación, la política \texttt{FORWARD} se restauró a \texttt{DROP} y el \textit{forwarding} se desactivó (\texttt{net.ipv4.ip\_forward=0}), devolviendo el segmento a su estado de aislamiento total.

\section{Rigor Documental}

Toda la documentación del proyecto (esta memoria, las decisiones de diseño, los diagramas) se redacta en \LaTeX{} y se aloja en el mismo repositorio Git que el código fuente. Este enfoque de \textit{Docs as Code} asegura trazabilidad entre las decisiones de diseño, el código que las implementa y la documentación que las describe, además de proporcionar consistencia tipográfica profesional y control de versiones unificado.

\section{Uso de la Inteligencia Artificial en el Desarrollo}

Durante el desarrollo de este proyecto se ha utilizado inteligencia artificial como herramienta de apoyo, tanto para la generación de fragmentos de código como para la redacción de partes de esta memoria, siempre bajo criterio de supervisión humana directa. Las herramientas empleadas incluyen:

\begin{itemize}
    \item \textbf{OpenCode}~\cite{opencode}: asistente de codificación interactivo ejecutado directamente en el servidor de desarrollo, utilizado para depuración, refactorización y generación de fragmentos de código Python y Bash, así como para la edición y corrección de la memoria en \LaTeX{}.
    \item \textbf{DeepSeek}~\cite{deepseek_api}: empleado como herramienta de consulta externa durante el desarrollo para documentación técnica y resolución de problemas. No forma parte del sistema en producción, que utiliza exclusivamente modelos locales a través de Ollama para garantizar la soberanía del dato.
    \item Otros modelos de lenguaje (ChatGPT, Gemini) utilizados puntualmente para consultas específicas durante la fase de investigación inicial.
\end{itemize}

Cada línea de código y cada párrafo generado han sido revisados, comprendidos, validados y, cuando ha sido necesario, modificados por el autor. No hay nada en este trabajo que no haya pasado por mi criterio personal.

Esta forma de trabajar no es distinta de la filosofía \textit{human in the loop} (el sistema propone, el humano dispone) que aplica el propio sistema Luxe Core AI en su módulo Smart Advisor (detallado en la Sección~\ref{sec:smart_advisor}): la IA sugiere, el humano decide. Del mismo modo que el Smart Advisor pregunta antes de actuar ante una subida de temperatura, aquí la inteligencia artificial ha actuado como asistente de propuestas, nunca como autoridad incuestionable.
